% Packages: longtable, array, booktabs. Unicode text requires XeLaTeX or LuaLaTeX.
% Attribution follows the supplied chapter; it does not establish discovery priority.
\begingroup
\small
\setlength{\tabcolsep}{3pt}
\renewcommand{\arraystretch}{1.15}
\begin{longtable}{@{}>{\raggedright\arraybackslash}p{0.045\linewidth}>{\raggedright\arraybackslash}p{0.085\linewidth}>{\raggedright\arraybackslash}p{0.39\linewidth}>{\raggedright\arraybackslash}p{0.135\linewidth}>{\raggedright\arraybackslash}p{0.285\linewidth}@{}}
\caption{Colophon formulae, constructions, and markers documented in the chapter, with scholarly attribution as stated there.}\label{tab:colophon-formulae-attribution}\\
\toprule
ID & Language & Formula or related forms & Function & Attribution in the chapter \\
\midrule
\endfirsthead
\multicolumn{5}{l}{\small\itshape Colophon expressions (continued)}\\
\toprule
ID & Language & Formula or related forms & Function & Attribution in the chapter \\
\midrule
\endhead
\midrule
\multicolumn{5}{r}{\small\itshape Continued on following page}\\
\endfoot
\bottomrule
\endlastfoot
L01 & Latin & \textit{Finito libro sit laus et gloria Christo\newline Finito libro laus sit gloria Christo\newline Finito libro laus sit Christo\newline Laus et gloria sit domino Iesu Christo} & Completion and praise\newline{\footnotesize Formula family} & Lucien Reynhout \\
\addlinespace
L02 & Latin & \textit{Finis adest operis; mercedem posco laboris} & Labour and reward\newline{\footnotesize Formula family} & Lucien Reynhout\newline Review: The spelling “merce dem poseo laboris” in two tables conflicts with the coherent reading in the prose and regional table. It is not treated as a manuscript variant. \\
\addlinespace
L03 & Latin & \textit{Explicit, expliceat; ludere scriptor eat\newline Explicit expliceat ludere scriptrix eat\newline Explicit expliceat. Scriptrix deo gratias} & Completion and wordplay\newline{\footnotesize Formula family} & Lucien Reynhout\newline Review: The feminine form is a corpus adaptation of the formula; it is not separately credited to Reynhout as a first identification. The Gertrude example preserves Explicit expliceat but not the ludere clause; it must not be counted as an exact full-formula match. \\
\addlinespace
L04 & Latin & \textit{Explicit iste liber; sit scriptor crimine liber} & Completion and absolution\newline{\footnotesize Formula family} & Lucien Reynhout \\
\addlinespace
L05 & Latin & \textit{Detur pro penna scriptori pulchra puella\newline Detur pro pena Egflib pulcra puella} & Scribal reward and wordplay\newline{\footnotesize Formula family} & Lucien Reynhout \\
\addlinespace
L06 & Latin & \textit{Finito libro referatur gratia Christo} & Completion and thanksgiving\newline{\footnotesize Formula family} & Lucien Reynhout \\
\addlinespace
L07 & Latin & \textit{Qui scripsit scribat, semper cum Domino vivat\newline Qui scripsit scribat semper cum Domino uiuat\newline Quod scripsit scribat, semper cum Domino vivat\newline Vivat in celis \ldots{} cum Domino felix} & Prayer for the scribe\newline{\footnotesize Formula family} & Lucien Reynhout ; Moreton discusses Italian nun-scribes\newline Review: The transmitted Quod form is preserved as quoted rather than silently corrected. The name-bearing celestial-life extension is grouped here. \\
\addlinespace
L08 & Latin & \textit{Cuius animae propitietur Deus} & Prayer for the soul\newline{\footnotesize Formula family} & Lucien Reynhout\newline Review: The chapter first attributes this formula to Reynhout and calls it absent from the corpus, then describes it among formulae not documented by Reynhout. Corpus status and that sentence require reconciliation. \\
\addlinespace
L09 & Latin & \textit{Nomen scriptoris \ldots{} plenus amoris} & Name formula\newline{\footnotesize Formula family} & Lucien Reynhout\newline Review: The same contradictory attribution/absence statements occur as for Cuius animae. The ellipsis is a variable portion, not an exact uninterrupted phrase. \\
\addlinespace
L10 & Latin & \textit{Finitus et completus} & Completion\newline{\footnotesize Formula family} & Lucien Reynhout \\
\addlinespace
L11 & Latin & \textit{Oretis pro scriptore propter Deum\newline Oretis pro scriptore} & Prayer request\newline{\footnotesize Formula family} & Lucien Reynhout\newline Review: The full formula is reported absent from the corpus. Short prayer expressions must not be counted as occurrences of this full formula. \\
\addlinespace
L12 & Latin & \textit{Feliciter} & Successful completion / benediction\newline{\footnotesize Formula family} & Lucien Reynhout\newline Review: The chapter identifies its clearest example as a male scribe in a manuscript associated with women. It must not automatically be labelled a female-scribe occurrence. \\
\addlinespace
L13 & Latin & \textit{Transcriptus} & Copying / completion\newline{\footnotesize Formula family} & Lucien Reynhout\newline Review: Reported absent from the corpus. The chapter supplies no feminine variant; no feminine form is inferred in this inventory. \\
\addlinespace
L14 & Latin & \textit{IHS\newline Iesus} & Invocation\newline{\footnotesize Devotional sign / short invocation} & Reynhout typological framework; examples supplied by the chapter\newline Review: Iesus occurs in the quoted Angela colophon. It is related invocation evidence, not an expansion certified for every IHS sign. \\
\addlinespace
L15 & Latin & \textit{In nomine Ihesu Christi} & Invocation\newline{\footnotesize Formula family} & Reynhout typological framework; examples supplied by the chapter \\
\addlinespace
L16 & Latin & \textit{Laus Deo\newline Sit laus Deo\newline Laus Deo sit semper\newline Laus Deo semper\newline Laus Deo in sempiternum\newline Sit laus Deo patri\newline Laus Deo sit semper benedictus\newline Laus deo et virgini gloriose} & Praise\newline{\footnotesize Formula family} & Reynhout typological framework; examples supplied by the chapter\newline Review: These are lexical or extended variants of the praise family; matching a short base expression does not establish any specific expanded variant. \\
\addlinespace
L17 & Latin & \textit{Deo gratias\newline Deo gracias\newline Mille Deo gratias\newline Deo gratias. Amen. Finis} & Thanksgiving\newline{\footnotesize Formula family} & Reynhout typological framework; examples supplied by the chapter\newline Review: Amen and Finis are adjoining elements in the last example, not evidence that every Deo gratias belongs to that extended sequence. \\
\addlinespace
L18 & Latin & \textit{Ad honorem; ad laudem; ad gloriam\newline Ad honorem et laudem Dei} & Praise and dedication\newline{\footnotesize Related constructions} & Reynhout typological framework; examples supplied by the chapter \\
\addlinespace
L19 & Latin & \textit{Finitus; finitus est} & Completion\newline{\footnotesize Completion marker} & Reynhout typological framework; examples supplied by the chapter \\
\addlinespace
L20 & Latin & \textit{Completus; completus est\newline Completus est liber praesens} & Completion\newline{\footnotesize Completion marker} & Reynhout typological framework; examples supplied by the chapter \\
\addlinespace
L21 & Latin & \textit{Explicit\newline Explicit liber iste\newline Explicit le Trespas de N. Dame\newline Explicit la benoites Passion de Nostre Seigneur Jhesu Crist\newline Explicit la espositione di sancto Bernardo sopra la Cantica canticorum di Salamone\newline Explicit Vita sancti Bernardi abbatis\newline Explicit \ldots{} commencié/perfet} & Textual closure\newline{\footnotesize Formula family} & Reynhout typological framework; examples supplied by the chapter\newline Review: The title-bearing French and Italian examples are hybrid contexts. A textual explicit is not necessarily a scribal production statement. \\
\addlinespace
L22 & Latin & \textit{Finis; est finis} & Closure\newline{\footnotesize Closure marker} & Unknown Hands chapter; no individual scholar assigned \\
\addlinespace
L23 & Latin & \textit{Anno Domini\newline Anno domini \ldots{} jar\newline Anno domino} & Dating\newline{\footnotesize Dating construction} & Reynhout typological framework; examples supplied by the chapter\newline Review: The latter form appears inside the Swedish colophon as quoted; it is not silently regularised. \\
\addlinespace
L24 & Latin & \textit{Orate pro scriptore / Orate pro scriptrice\newline Oretur igitur pro scriptrice\newline Orate pro scriptrice unum paternoster\newline \ldots{} orent pro scriptrice Agnes et notatrice Christina} & Prayer for the scribe\newline{\footnotesize Formula family} & Reynhout typological framework; examples supplied by the chapter\newline Review: Distinct grammatical constructions are grouped as a prayer-request family. The scribe-and-notator extension remains distinguishable. \\
\addlinespace
L25 & Latin & \textit{Oretis pro ea ex caritate Dei} & Prayer request\newline{\footnotesize Formula family} & Unknown Hands chapter; no individual scholar assigned \\
\addlinespace
L26 & Latin & \textit{Orate per me} & Prayer request\newline{\footnotesize Attested expression} & Unknown Hands chapter; no individual scholar assigned\newline Review: Recorded wording is per me in the quoted St Gall colophon. It must not be normalised silently to pro me. \\
\addlinespace
L27 & Latin & \textit{Si quis abstulerit, anathema sit} & Protection against removal\newline{\footnotesize Formula family} & Reynhout typological framework; examples supplied by the chapter \\
\addlinespace
L28 & Latin & \textit{Qui te sustulerit \ldots{} hunc anathema ferit} & Protection against removal\newline{\footnotesize Attested verse expression} & Unknown Hands chapter; no individual scholar assigned \\
\addlinespace
L29 & Latin & \textit{Qui secus fecerit anathema coram summo Deo sit} & Sanction\newline{\footnotesize Attested expression} & Unknown Hands chapter; no individual scholar assigned \\
\addlinespace
L30 & Latin & \textit{Per me} & Responsibility / self-identification\newline{\footnotesize Responsibility construction} & Reynhout typological framework; examples supplied by the chapter \\
\addlinespace
L31 & Latin & \textit{Scripsit; scriptus; scriptus per manum\newline Manibus suis scripsit\newline Johannes me scripsit} & Writing and responsibility\newline{\footnotesize Related responsibility constructions} & Reynhout typological framework; examples supplied by the chapter\newline Review: Johannes me scripsit is quoted in the Gameson reading notes (p. 18); it is an example, not a formula whose first discovery is attributed to him. \\
\addlinespace
L32 & Latin & \textit{Hic liber est \ldots{} / Iste liber est \ldots{}} & Ownership\newline{\footnotesize Ownership construction} & Reynhout typological framework; examples supplied by the chapter \\
\addlinespace
L33 & Latin & \textit{Pertinet ad \ldots{} / Iste liber pertinet ad \ldots{}} & Ownership\newline{\footnotesize Ownership construction} & Reynhout typological framework; examples supplied by the chapter \\
\addlinespace
L34 & Latin & \textit{Lege feliciter} & Address to the reader\newline{\footnotesize Formula family} & Reynhout typological framework; examples supplied by the chapter\newline Review: Comparison expression; the chapter states that no clear corpus equivalent is tracked. \\
\addlinespace
L35 & Latin & \textit{Utere felix} & Address to the reader\newline{\footnotesize Formula family} & Reynhout typological framework; examples supplied by the chapter\newline Review: Comparison expression; the chapter states that no clear corpus equivalent is tracked. \\
\addlinespace
L36 & Latin & \textit{Bene valeas qui legis} & Address to the reader\newline{\footnotesize Formula family} & Reynhout typological framework; examples supplied by the chapter\newline Review: Comparison expression; the chapter states that no clear corpus equivalent is tracked. \\
\addlinespace
L37 & Latin & \textit{Que / Qui me scribebat \ldots{} nomen habebat\newline Que me scribebat. Gertrudis nomen habebat} & Naming the scribe\newline{\footnotesize Formula family} & Unresolved: chapter prose versus final comparison table\newline Review: The prose calls these formulae not documented by Reynhout; the final table marks them as Reynhout. Preserve the name slot and flag the attribution. \\
\addlinespace
L38 & Latin & \textit{Finitus est iste liber per me soror \ldots{}} & Completion and female self-identification\newline{\footnotesize Formula family} & Unresolved: chapter prose versus final comparison table\newline Review: The prose presents this among formulae not documented by Reynhout; the final table credits Reynhout. \\
\addlinespace
L39 & Latin & \textit{Iste liber scripsit \ldots{}} & Responsibility\newline{\footnotesize Formula family} & Unresolved: final table credits Reynhout without formula-specific support\newline Review: Quoted for Angela Rucellai in the prose. The chapter does not establish a formula-specific Reynhout attribution. \\
\addlinespace
L40 & Latin & \textit{Indigna monialis; indignus servus; serva Dei indigna\newline Soror [Name] indigna serva Domini nostri Iesu Christi} & Humility\newline{\footnotesize Related humility constructions} & Reynhout vocabulary comparison; chapter examples \\
\addlinespace
L41 & Latin & \textit{Peccatrix\newline Soror [Name] peccatrix} & Humility / penitence\newline{\footnotesize Lexical marker} & Unknown Hands chapter; no individual scholar assigned \\
\addlinespace
L42 & Latin & \textit{Non bene finivi, quia scribere nescivi} & Humility and scribal ability\newline{\footnotesize Attested expression} & Unknown Hands chapter; no individual scholar assigned \\
\addlinespace
L43 & Latin & \textit{Christi sit dilectrix [Name], libri huius scriptrix} & Female scribal self-identification\newline{\footnotesize Attested expression} & Unknown Hands chapter; no individual scholar assigned \\
\addlinespace
L44 & Latin & \textit{Ego soror [Name] \ldots{} scripsi manu propria / composui hunc librum} & Responsibility\newline{\footnotesize Responsibility construction} & Chapter examples; Moreton comparison in working notes\newline Review: The chapter supplies both constructions. The Angela example is associated with different surnames in different passages; witness attribution requires review. \\
\addlinespace
L45 & Latin & \textit{Benedicamus Domino\newline Benedicamus Domino. Deo gratias} & Praise / liturgical response\newline{\footnotesize Attested expression} & Unknown Hands chapter; no individual scholar assigned \\
\addlinespace
L46 & Latin & \textit{Ave Maria; Paternoster; Requiem} & Prayer requested for remembrance\newline{\footnotesize Prayer names} & Unknown Hands chapter; no individual scholar assigned\newline Review: Names of prayers embedded in requests, not three independent colophon formulae. A bare prayer name is insufficient to classify the whole request. \\
\addlinespace
L47 & Latin & \textit{Qui legentium orationes pro merito caritatis expostulat} & Prayer and remembrance\newline{\footnotesize Formula family} & Melissa Moreton and Kathleen Arthur (chapter attribution)\newline Review: Appears in the appended working discussion, alongside Qui scripsit scribat. The chapter does not provide a formula-specific citation for Arthur. \\
\addlinespace
L48 & Latin & \textit{Quicunque cantat vel legat in eo habeat nostri memoriam aput Deum} & Reader remembrance\newline{\footnotesize Attested expression} & Unknown Hands chapter; no individual scholar assigned\newline Review: Quoted for Katharina von Brugg; a specific attestation rather than a demonstrated independently recurring formula. \\
\addlinespace
L49 & Latin & \textit{Deus sit laudatus et pro me deprecatus} & Praise and prayer request\newline{\footnotesize Attested expression} & Unknown Hands chapter; no individual scholar assigned\newline Review: Quoted within the Angela colophon; its recurrence is not demonstrated by this chapter. \\
\addlinespace
D01 & Dutch & \textit{Int jaer ons heren\newline Jnt iaer ons heren\newline In het jaar des Heren} & Dating\newline{\footnotesize Dating construction} & Unknown Hands chapter; no individual scholar assigned \\
\addlinespace
D02 & Dutch & \textit{Bidt om gods wil / Om gods wil \ldots{}\newline Om gods wil eenen aue maria voer diet ghescreuen heeft} & Prayer request\newline{\footnotesize Formula family} & Unknown Hands chapter; no individual scholar assigned\newline Review: The phrase om gods wil alone does not establish a prayer request without its context. \\
\addlinespace
D03 & Dutch & \textit{Dit boec / boeck hoert toe \ldots{}\newline Dit boec hoert toe den convente\newline Dit boeck hoert toe den cloester\newline Dit boec hoert toe den cloester van Jhericho} & Ownership\newline{\footnotesize Ownership construction} & Unknown Hands chapter; no individual scholar assigned \\
\addlinespace
D04 & Dutch & \textit{Dit boec hoert tot \ldots{}} & Ownership / community affiliation\newline{\footnotesize Ownership construction} & Unknown Hands chapter; no individual scholar assigned \\
\addlinespace
D05 & Dutch & \textit{Dit boec behoort toe \ldots{}} & Ownership\newline{\footnotesize Ownership construction} & Unknown Hands chapter; no individual scholar assigned \\
\addlinespace
D06 & Dutch & \textit{Dit boec is \ldots{}} & Ownership\newline{\footnotesize Ownership construction} & Unknown Hands chapter; no individual scholar assigned \\
\addlinespace
D07 & Dutch & \textit{Dit boeck heeft ghescreven suster [Name] \ldots{}} & Responsibility\newline{\footnotesize Attested construction} & Unknown Hands chapter; no individual scholar assigned\newline Review: Extracted from the Anna Joerdaens example, rather than attributed as a separate recurrent type. \\
\addlinespace
G01 & German & \textit{Pitt got für [sie]\newline Bitten Gott für sy\newline Bittend got für \ldots{}} & Prayer request\newline{\footnotesize Formula family} & Unknown Hands chapter; no individual scholar assigned \\
\addlinespace
G02 & German & \textit{Bidt vor die schrivers\newline Bidt vor die schrivers dattem got ewich leven geve} & Collective prayer request\newline{\footnotesize Formula family} & Unknown Hands chapter; no individual scholar assigned\newline Review: Labelled German by the chapter. Its Low German / Middle Dutch linguistic placement requires review. \\
\addlinespace
G03 & German & \textit{Das puch hat geschriben swester \ldots{}} & Responsibility\newline{\footnotesize Responsibility construction} & Unknown Hands chapter; no individual scholar assigned \\
\addlinespace
G04 & German & \textit{Von Schwester [Name]\newline Von swester [Name]} & Responsibility\newline{\footnotesize Responsibility construction} & Unknown Hands chapter; no individual scholar assigned \\
\addlinespace
G05 & German & \textit{Die schreiberin die [Text] geschriben hat} & Responsibility\newline{\footnotesize Responsibility construction} & Unknown Hands chapter; no individual scholar assigned \\
\addlinespace
G06 & German & \textit{Do man zalt / Als man zalt} & Dating\newline{\footnotesize Dating constructions} & Unknown Hands chapter; no individual scholar assigned \\
\addlinespace
G07 & German & \textit{Nach Christi gepurt\newline Von unsers lieben herren Jhesu Cristi gepurd} & Dating\newline{\footnotesize Dating construction} & Unknown Hands chapter; no individual scholar assigned \\
\addlinespace
G08 & German & \textit{Vollendet; volbracht; geendet\newline Dis bůch ist volendet und geschriben\newline Ist disz puch auszgeschriben worden} & Completion\newline{\footnotesize Completion markers} & Unknown Hands chapter; no individual scholar assigned \\
\addlinespace
G09 & German & \textit{Lob sye Gott\newline Got zu lob} & Praise\newline{\footnotesize Formula family} & Unknown Hands chapter; no individual scholar assigned \\
\addlinespace
G10 & German & \textit{Zu lob vnd erenn} & Praise / dedication\newline{\footnotesize Formula family} & Unknown Hands chapter; no individual scholar assigned \\
\addlinespace
G11 & German & \textit{Dis(z) / Das Buch gehört \ldots{}\newline Das puch gehört\newline Dis puch gehört\newline Dasz puch gehört\newline Gehört in das Kloster\newline Gehört in die Sacristei\newline Gehört zu \ldots{}} & Ownership\newline{\footnotesize Ownership constructions} & Unknown Hands chapter; no individual scholar assigned \\
\addlinespace
G12 & German & \textit{Dis Buch ist \ldots{}} & Ownership\newline{\footnotesize Ownership construction} & Unknown Hands chapter; no individual scholar assigned \\
\addlinespace
G13 & German & \textit{Bosse Schriberin\newline Bittend got für ich bosse Schriberin} & Humility\newline{\footnotesize Lexical marker / attested expression} & Unknown Hands chapter; no individual scholar assigned \\
\addlinespace
G14 & German & \textit{O wie fro ich was, da ich schrib Deo gratias} & Joy at completion\newline{\footnotesize Attested expression} & Unknown Hands chapter; no individual scholar assigned\newline Review: A quoted mixed German–Latin expression, not demonstrated as an independent recurring formula. \\
\addlinespace
I01 & Italian & \textit{Finito\newline Finito le hepistole del seguente libro\newline Finito è il presente libro} & Completion\newline{\footnotesize Completion marker} & Melissa Moreton \\
\addlinespace
I02 & Italian & \textit{Finisce / Qui finisce\newline Qui finisce il libro\newline Finisse di scrivere\newline Finissi di scrivere} & Completion\newline{\footnotesize Completion construction} & Melissa Moreton \\
\addlinespace
I03 & Italian & \textit{Si compié di scrivere \ldots{}} & Completion\newline{\footnotesize Attested construction} & Unknown Hands chapter; no individual scholar assigned \\
\addlinespace
I04 & Italian & \textit{Indegniamente / indegna\newline Una suora indegniamente monacha} & Humility\newline{\footnotesize Lexical markers} & Melissa Moreton \\
\addlinespace
I05 & Italian & \textit{Peccatrice / peccatore} & Humility and penitence\newline{\footnotesize Lexical markers} & Melissa Moreton ; chapter example for peccatore \\
\addlinespace
I06 & Italian & \textit{A llaude et onore\newline Ad honore et laude di Dio\newline A llaude et gloria della somma et eterna sapientia} & Praise / dedication\newline{\footnotesize Related praise constructions} & Unknown Hands chapter; no individual scholar assigned \\
\addlinespace
I07 & Italian & \textit{Questo libro è delle monache / delle suore \ldots{}\newline Questo libro è del convento \ldots{}\newline Questo libro è del monastero \ldots{}} & Institutional ownership\newline{\footnotesize Ownership construction} & Unknown Hands chapter; no individual scholar assigned \\
\addlinespace
I08 & Italian & \textit{Questo libro è di \ldots{}} & Personal ownership\newline{\footnotesize Ownership construction} & Unknown Hands chapter; no individual scholar assigned \\
\addlinespace
I09 & Italian & \textit{Questo libro è di \ldots{} e del convento \ldots{}} & Shared ownership\newline{\footnotesize Ownership construction} & Unknown Hands chapter; no individual scholar assigned \\
\addlinespace
I10 & Italian & \textit{Chi l’acchata lo renda\newline Chi llo achatta lo riguardi et rendalo \ldots{}\newline Rendilo al monasterio del Paradiso} & Return of the book\newline{\footnotesize Return construction} & Unknown Hands chapter; no individual scholar assigned \\
\addlinespace
I11 & Italian & \textit{Scripto per mano di suora [Name]\newline Ho scritto io suora [Name]} & Responsibility\newline{\footnotesize Responsibility construction} & Unknown Hands chapter; no individual scholar assigned \\
\addlinespace
I12 & Italian & \textit{Priega per me misera, peccatore} & Prayer and humility\newline{\footnotesize Attested expression} & Unknown Hands chapter; no individual scholar assigned\newline Review: The earlier extended Raffaella quotation refers to ms16292; the humility discussion gives ms16297. The manuscript reference requires review. \\
\addlinespace
I13 & Italian & \textit{Indegna serva e schiava di Jesu cristo} & Humility and Catherine of Siena imitation\newline{\footnotesize Formula family} & Melissa Moreton\newline Review: Working-note quotation attributes the discussion to Moreton, p. 217. Its bibliographical year and edited-volume reference require reconciliation with the chapter citation keys. \\
\addlinespace
F01 & French & \textit{Pries Nostre Seigneur} & Prayer request\newline{\footnotesize Formula family} & Unknown Hands chapter; no individual scholar assigned \\
\addlinespace
F02 & French & \textit{Pour ses soeurs} & Communal dedication / intended audience\newline{\footnotesize Dedication construction} & Unknown Hands chapter; no individual scholar assigned\newline Review: The Catherine de Ligne quotation states that the breviary was written for her sisters. It does not by itself establish an intercessory prayer formula. \\
\addlinespace
F03 & French & \textit{Est à moy} & Personal ownership\newline{\footnotesize Ownership construction} & Unknown Hands chapter; no individual scholar assigned \\
\addlinespace
F04 & French & \textit{L’an de grace \ldots{}} & Dating\newline{\footnotesize Attested construction} & Unknown Hands chapter; no individual scholar assigned\newline Review: Supplied by the Bernardine Dunant quotation; not separately presented as a recurrent formula in the chapter. \\
\addlinespace
P01 & Portuguese & \textit{Acabousse} & Completion\newline{\footnotesize Completion marker} & Unknown Hands chapter; no individual scholar assigned \\
\addlinespace
P02 & Portuguese & \textit{Screveo freira \ldots{}} & Responsibility\newline{\footnotesize Responsibility construction} & Unknown Hands chapter; no individual scholar assigned\newline Review: Named in the final comparison table; no complete Portuguese witness is supplied for this expression. \\
\addlinespace
P03 & Portuguese & \textit{Este livro he do mosteiro \ldots{}} & Institutional ownership\newline{\footnotesize Ownership construction} & Unknown Hands chapter; no individual scholar assigned \\
\addlinespace
S01 & Swedish & \textit{Conuentz syster\newline Owerdogh conuentz syster} & Convent identity\newline{\footnotesize Identity marker} & Chapter; Moreton discussion citing Lindell for the Brask example\newline Review: Lindell supplies a transcription and translation (2010, pp. 181–182); the chapter does not establish priority of identifying a formula. \\
\addlinespace
S02 & Swedish & \textit{Owärdoghe\newline Owerdogh} & Humility\newline{\footnotesize Lexical marker} & Chapter; Moreton discussion citing Lindell for a related form \\
\addlinespace
S03 & Swedish & \textit{Bidhin kära systra} & Prayer request to sisters\newline{\footnotesize Formula family} & Unknown Hands chapter; no individual scholar assigned \\
\addlinespace
S04 & Swedish & \textit{\ldots{} hore til \ldots{}} & Ownership\newline{\footnotesize Ownership construction} & Unknown Hands chapter; no individual scholar assigned \\
\addlinespace
E01 & Old English & \textit{Wulfwi me wrat} & Responsibility\newline{\footnotesize Attested expression} & Richard Gameson, as quoted in the reading notes (p. 18)\newline Review: An individual attestation, not a formula whose first identification is demonstrated in the chapter. \\
\addlinespace
N01 & Non-verbal & \textit{Crosses; chrismons; Marian initials} & Invocation\newline{\footnotesize Sign categories} & Reynhout typological framework; examples supplied by the chapter\newline Review: The chapter supplies sign categories rather than transcribed symbols. These cannot be added to a text-only matcher as if they were strings. \\
\addlinespace
\end{longtable}
\providecommand{\FormulaChapterCite}[1]{\cite{#1}}
\noindent\textit{Sources cited by the chapter.} Reynhout: \FormulaChapterCite{reynhout_codicologie_2001}, \FormulaChapterCite{reynhout_formules_2006}, \FormulaChapterCite{reynhout_colophons_2025}\\
Moreton: \FormulaChapterCite{kiraz_nun-scribes_2024}; Gameson: \FormulaChapterCite{gameson_scribe_2002}.\\
The working notes cite Lindell (2010, pp. 181--182) for the Brask transcription. The Moreton edited-volume reference and Arthur attribution require bibliographical review.
\par\medskip
\noindent\textit{Editorial scope.} A row can represent a formula family, a construction, a lexical marker, an attested expression, or a non-verbal sign category. Several forms in one row are not necessarily interchangeable spellings. The catalogue distinguishes these types. Language labels follow the chapter, with disputed placements flagged. ``Unknown Hands chapter'' indicates an observation in the supplied chapter without a formula-specific named-scholar attribution. No occurrence counts are transcribed from the working chapter.
\endgroup
